\chapter{单元系的相变与复相平衡}
\problem{7.1}{宇宙学观测表明，我们所处的宇宙在加速膨胀，这表明，宇宙的组分除包含物质外，还包含被称为暗能量的成分。假设暗能量的状态方程可以表示为 $p_{\rm DE}=\kappa_{\rm DE}\varepsilon$，其中 $\varepsilon$ 为平均能量密度，根据相（不）稳定条件确定系数 $\kappa_{\rm DE}$ 的取值范围。}
\answerarea
\problem{7.2}{试对热力学系统中的自发过程，证明下列判据：
\begin{enumerate}
\item 如果熵和体积保持不变，则使内能取得极小值的状态为系统的稳定平衡态。
\item 如果熵和压强保持不变，则使焓取得极小值的状态为系统的稳定平衡态。
\item 如果焓和压强保持不变，则使熵取得极大值的状态为系统的稳定平衡态。
\item 如果自由能和体积保持不变，则使温度最低的状态为系统的稳定平衡态。
\item 如果自由焓和压强保持不变，则使温度最低的状态为系统的稳定平衡态。
\item 如果自由能和温度保持不变，则使体积最小的状态为系统的稳定平衡态。
\item 如果内能和熵保持不变，则使体积最小的状态为系统的稳定平衡态。
\end{enumerate}}
\answerarea
\clearpage
\problem{7.3}{对于在温度 $T$、压强 $p$ 的情况下发生的相变，试证明系统的摩尔内能改变为
\[
\Delta U_{\rm m}=L_{\rm m}\left(1-\frac pT\frac{\mathrm dT}{\mathrm dp}\right),
\]
其中 $L_{\rm m}$ 为摩尔汽化热。}
\answerarea
\problem{7.4}{在 $p_0=101\,325\,\mathrm{Pa}$ 的压强下，冰的熔点为 $273.15\,\mathrm K$，此时冰的熔解热为 $l=3.35\times10^5\,\mathrm{J/kg}$，冰的比体积为 $v_{\rm 固}=1.0907\times10^{-3}\,\mathrm{m^3/kg}$，水的比体积为 $v_{\rm 液}=1.00013\times10^{-3}\,\mathrm{m^3/kg}$。试确定冰水两相平衡曲线的斜率 $\mathrm dp/\mathrm dT$。}
\answerarea
\clearpage
\problem{7.5}{在 $p_0=101\,325\,\mathrm{Pa}$ 的压强下，水的沸点为 $99.974\,^\circ\mathrm C$，此时水的汽化热为 $l=2.257\times10^6\,\mathrm{J/kg}$，水的比体积为 $v_{\rm 液}=1.043\times10^{-3}\,\mathrm{m^3/kg}$，水蒸气的比体积为 $v_{\rm 气}=1.673\,\mathrm{m^3/kg}$，试确定水和水蒸气两相平衡的情况下水的沸点随压强的变化率 $\mathrm dT/\mathrm dp$。}
\answerarea
\problem{7.6}{在 $1\,\mathrm{atm}$ 下，固态氩于 $83\,\mathrm K$ 熔解，其摩尔熔解热为 $1176\,\mathrm{J/mol}$，全部熔解时其摩尔体积增加了 $3.5\times10^{-6}\,\mathrm{m^3}$。有外力施加时，其潜热保持不变，但体积的变化与 $T^{3/2}$ 成反比。试确定：若要使它的熔点提高一倍，需施加多大压力？}
\answerarea
\clearpage
\problem{7.7}{液氦的正常沸点为 $4.2\,\mathrm K$，压强降为 $1\,\mathrm{Torr}$ 时的沸点为 $1.2\,\mathrm K$，试估计 $^4\mathrm{He}$ 在这一温度范围内的平均汽化热。}
\answerarea
\problem{7.8}{设某种液体在某一压强下的沸点为 $400\,\mathrm K$，且沸点每升高 $1\,\mathrm K$，其平衡压强升高 $5\%$，试估计这种液体的汽化热。}
\answerarea
\clearpage
\problem{7.9}{假定水蒸气可视作理想气体，试由下表所列数值计算 $27\,^\circ\mathrm C$ 时水的升华热。
\begin{center}
\begin{tabular}{|c|c|c|c|}
\hline
温度（$^\circ\mathrm C$）&27.1&27.0&26.9\\
\hline
饱和蒸气压（$10^3\,\mathrm{Pa}$）&3.587&3.566&3.545\\
\hline
\end{tabular}
\end{center}}
\answerarea
\problem{7.10}{在三相点附近，固态氨的饱和蒸气压（单位为 atm）方程为：$\ln p=23.03-\dfrac{3754}{T}$，液态氨的饱和蒸气压方程为 $p=19.49-\dfrac{3063}{T}$，试求氨的三相点温度、压强以及氨在三相点的汽化热、升华热和熔解热。}
\answerarea
\clearpage
\problem{7.11}{某物质的摩尔质量为 $\mu$，三相点温度及压强分别为 $T_0,p_0$，三相点时固态及液态的密度分别为 $\rho_s,\rho_l$，其蒸气可视为理想气体，又在三相点时的熔解曲线的斜率为 $\left(\dfrac{\mathrm dp}{\mathrm dT}\right)_{\rm sol}$，饱和蒸气压曲线斜率为 $\left(\dfrac{\mathrm dp}{\mathrm dT}\right)_{\rm sv}$。
\begin{enumerate}
\item 试确定升华曲线的斜率 $\left(\dfrac{\mathrm dp}{\mathrm dT}\right)_{\rm sub}$；
\item 证明在通常情况下有：$\left(\dfrac{\mathrm dp}{\mathrm dT}\right)_{\rm sub}>\left(\dfrac{\mathrm dp}{\mathrm dT}\right)_{\rm sv}$。
\end{enumerate}}
\answerarea
\problem{7.12}{已知池中水面上冻结了 $1\,\mathrm{cm}$ 厚的冰层，冰上面空气的温度为 $-20\,^\circ\mathrm C$，冰的热导率 $\kappa=2.092\,\mathrm{J/(m\cdot s\cdot K)}$，结冰时的潜热 $l=3.349\times10^5\,\mathrm{J/kg}$，冰的密度 $\rho=10^3\,\mathrm{kg/m^3}$，试确定：
\begin{enumerate}
\item 开始时，冰层厚度增加的速度；
\item 冰层厚度增加一倍所需的时间。
\end{enumerate}}
\answerarea
\clearpage
\problem{7.13}{质量为 $M$ 的固体，密度为 $\rho_1$，在温度为 $T$、压强为 $p_1$ 的情况下熔解为密度为 $\rho_2$ 的液体，熔解热为 $l$。试确定：
\begin{enumerate}
\item 固体熔解成液体时内能的增量；
\item 熵的增量。
\end{enumerate}}
\answerarea
\problem{7.14}{对于液气相变，试证明，当蒸气和液相在温度为 $T$ 情况下达到平衡时，蒸气的体膨胀系数为
\[
\frac1{V_{\rm m}}\frac{\mathrm dV_{\rm m}}{\mathrm dT}
=\frac1T\left(1-\frac{L_{\rm m}}{RT}\right),
\]
其中 $L_{\rm m}$ 为系统的摩尔汽化热，$V_{\rm m}$ 为摩尔体积，$R$ 为普适气体常量。}
\answerarea
\clearpage
\problem{7.15}{试证明：将 $p$-$V$ 图上不同温度下的范德瓦尔斯模型的等温线的极小值点和极大值点连起来的曲线的曲线方程为
\[
pV_{\rm m}^3=a(V_{\rm m}-2b).
\]}
\answerarea
\clearpage
\problem{7.16}{目前，由于精密焊接等技术的需求，人们对亚微米、甚至纳米尺度的缝隙或细管中的液体的性质（例如，凝固、蒸发等）极为关注。假设碳纳米管中水与管壁间的接触角为 $\theta$，纳米管的半径为 $R_p$，水的表面张力系数为 $\sigma$，表面平直的水的凝结温度为 $T_{\rm s,B}$，冰的熔解热为 $l$、密度为 $\rho$。
\begin{enumerate}
\item 试证明：上述碳纳米管中的水的凝结温度相对于表面平直的水的凝结温度的改变量为
\[
\Delta T_{\rm s}=T_{\rm s,p}-T_{\rm s,B}
=-\frac{2\sigma\cos\theta}{l\rho R_p}T_{\rm s,B}.
\]
\item 实验测得，在接触角 $\theta=83.5^\circ$ 的情况下，纳米管中水的凝结温度的改变量与纳米管内（直）径 $D_p$ 间的对应关系如下表所示。试总结出实验测得的 $\Delta T_{\rm s}$ 与 $D_p$ 间的关系，并说明实验测量结果与前述简单理论关系的异同。
\begin{center}\small
\begin{tabular}{c|rrrrrrr}
$D_p/\mathrm{nm}$&1.35&1.70&1.87&3.21&6.57&8.17&8.86\\
\hline
$\Delta T_{\rm s}/\mathrm K$&$-21.0$&$-15.0$&$-14.2$&$-10.5$&$-7.5$&$-6.5$&$-4.5$
\end{tabular}
\end{center}
\item 若假设上述实验结果与理论结果之间的差异仅仅来自于前述近似理论忽略了水的表面张力系数对纳米管内径的依赖性，那么这些实验结果表明纳米管中水的表面张力系数与纳米管内径间的关系是什么？（定性给出 $\sigma$ 与 $1/D_p$ 间的函数关系，进而予以讨论即可）
\end{enumerate}}
\answerarea
\clearpage
\problem{7.17}{与在二级相变中一样，也可以引入序参量描述一级相变。例如，对液气相变，这样的序参量可以取为摩尔体积与一很高温度下理想气体的摩尔体积的差值与该高温下理想气体的摩尔体积的比值。试在朗道相变理论框架下，说明由高对称相向低对称相的一级相变过程序参量随控制参量变化的行为，并给出图示。}
\answerarea
\problem{7.18}{对于一液滴，人们引入一个类似于偏心率的量 $\varepsilon$ 描述其形状，该量常被称为形变参数。通常，对于一个球形水滴在某一方向加压时，该方向向内收缩，与之垂直方向则向外凸出，这种模式的运动称为四极振动。显然，四极振动时，液滴的形变参数保持非零，即没有稳定的形变。当使该液滴绕某一轴转动时，如果角动量大于一个临界值，则液滴可以稳定地呈沿与转轴垂直的方向拉长的椭球形（$\varepsilon>0$）；如果角动量小于该临界值，则液滴仍然以四极振动为其主要集体运动模式。这种现象称为角动量驱动的形状相变。因发生相变时，形变参数 $\varepsilon$ 发生突变，因此这种相变为一级相变。试仿照朗道相变理论，构建一个模型描述角动量驱动的形状相变，并画出角动量和形变参数张成的空间中的相图。}
\answerarea
\clearpage
\problem{7.19}{熵是物理系统的重要的态函数，是系统的有序无序程度的度量；黏度是表征物理系统的输运性质的重要物理量。黏度与熵（或熵密度）的比值被视为描述系统由组分粒子间耦合较强的相到组分粒子间耦合很弱的相的相变（例如，液气相变）的重要物理量而备受关注。现有一系统，其温度由低于沸点（记为 $T_{\rm B}$）的 $T_{\rm L}$ 上升到高于沸点的 $T_{\rm G}$，系统发生明显的液气相变。试定量分析该相变过程中系统的黏度与摩尔熵密度的比值随温度变化的行为，并给出图示。}
\answerarea
\problem{7.20}{一系统在高温低密（低化学势）情况下可以发生连续相变，在低温高密（高化学势）情况下可以发生一级相变，相图中一级相变和连续过渡的分界点称为临界终点（相应的状态称为临界终点态）。如果上述连续过渡可以近似为二级相变，临界终点态的状态（控制）参量可以由温度和化学势表述，即 $(T_E,\mu_E)$，试仿照朗道二级相变理论构建一个模型，描述上述不同温度及不同化学势区域中发生不同级次相变的现象。}
\answerarea
\clearpage
